\documentclass{article}
\usepackage{amsmath,amssymb,amsthm}
\usepackage{algorithm}
\usepackage{algorithmic}
\usepackage{enumitem}
\usepackage{hyperref}
\usepackage{geometry}
\geometry{margin=1in}
\newtheorem{theorem}{Theorem}
\newtheorem{lemma}{Lemma}
\newcommand{\EE}{\mathbb{E}}
\newcommand{\RR}{\mathbb{R}}

\begin{document}

%%%%%%%%%%%%%%%%%%%%%%%%%%%%%%%%%%%%%%%%%%%%%%%%%%%%%%%%%%%%
\section*{AdamW (Adam with Decoupled Weight Decay)}
%%%%%%%%%%%%%%%%%%%%%%%%%%%%%%%%%%%%%%%%%%%%%%%%%%%%%%%%%%%%

\section*{Mathematical Formulation}
Let $f:\RR^{d}\rightarrow\RR$ be the (possibly non‑convex) loss function.  
At iteration $t\ge 1$ we keep:

\begin{itemize}
    \item Parameter vector $w_t\in\RR^{d}$,
    \item First‑moment estimate $m_t\in\RR^{d}$,
    \item Second‑moment estimate $v_t\in\RR^{d}$.
\end{itemize}

Given a stochastic gradient $g_t:=\nabla f(w_t;\xi_t)$ computed on a random mini‑batch $\xi_t$, the AdamW updates are:

\[
\begin{aligned}
    &\text{(Weight decay)} && w_t^{\text{dec}} = w_t - \eta_t \,\lambda\, w_t, \tag{1}\\[4pt]
    &\text{(First moment)} && m_t = \beta_1 m_{t-1} + (1-\beta_1) g_t, \tag{2}\\[4pt]
    &\text{(Second moment)}&& v_t = \beta_2 v_{t-1} + (1-\beta_2) g_t\odot g_t, \tag{3}\\[4pt]
    &\text{(Bias correction)}&& 
    \hat m_t = \frac{m_t}{1-\beta_1^{\,t}},\qquad 
    \hat v_t = \frac{v_t}{1-\beta_2^{\,t}},\tag{4}\\[4pt]
    &\text{(Parameter update)}&& 
    w_{t+1}= w_t^{\text{dec}} - \eta_t \,\frac{\hat m_t}{\sqrt{\hat v_t}+\epsilon}, \tag{5}
\end{aligned}
\]
where $\odot$ denotes element‑wise product, $\beta_1,\beta_2\in[0,1)$ are exponential decay rates, $\epsilon>0$ stabilises the division, $\lambda\ge0$ is the weight‑decay coefficient and $\{\eta_t\}$ is a stepsize schedule (often constant $\eta_t=\eta$).

\section*{Pseudocode}
\begin{algorithm}[H]
\caption{AdamW}
\begin{algorithmic}[1]
\STATE \textbf{Input:} initial $w_0$, stepsize $\eta$, decay rates $\beta_1,\beta_2$, weight‑decay $\lambda$, $\epsilon>0$, number of iterations $T$.
\STATE Initialize $m_0\gets 0$, $v_0\gets 0$.
\FOR{$t=1$ \TO $T$}
    \STATE Sample mini‑batch $\xi_t$ and compute stochastic gradient $g_t=\nabla f(w_{t-1};\xi_t)$.
    \STATE \textbf{Weight decay:} $w^{\text{dec}}_{t-1}\gets w_{t-1}-\eta\,\lambda\, w_{t-1}$.
    \STATE $m_t \gets \beta_1 m_{t-1} + (1-\beta_1) g_t$.
    \STATE $v_t \gets \beta_2 v_{t-1} + (1-\beta_2) g_t\odot g_t$.
    \STATE $\hat m_t \gets m_t/(1-\beta_1^{\,t})$, $\hat v_t \gets v_t/(1-\beta_2^{\,t})$.
    \STATE $w_t \gets w^{\text{dec}}_{t-1} - \eta\,\frac{\hat m_t}{\sqrt{\hat v_t}+\epsilon}$.
\ENDFOR
\STATE \textbf{Output:} $w_T$ (or the averaged iterate).
\end{algorithmic}
\end{algorithm}

\section*{Dry Run Example}
Consider the one‑dimensional quadratic
\[
f(w) = (w-3)^2,
\]
so $\nabla f(w)=2(w-3)$.  
Hyper‑parameters: $\eta=0.1$, $\beta_1=0.9$, $\beta_2=0.999$, $\epsilon=10^{-8}$, $\lambda=0.01$, $w_0=0$, $m_0=v_0=0$.  
We perform three iterations.

\begin{center}
\begin{tabular}{c|c|c|c|c|c}
$t$ & $w_{t-1}$ & $g_t$ & $m_t$ & $v_t$ & $w_t$ \\ \hline
1 & $0$ & $2(0-3)=-6$ & $0.1(-6)=-0.6$ & $0.001\cdot36=0.036$ & $w_1 = 0 -0.1\lambda\cdot0 -0.1\frac{-0.6}{\sqrt{0.036}+10^{-8}}\approx 1.0$ \\ \hline
2 & $1.0$ & $2(1-3)=-4$ & $0.9(-0.6)+0.1(-4)=-0.94$ & $0.999\cdot0.036+0.001\cdot16=0.051964$ & $w_2\approx 1.0-0.1\lambda(1.0)-0.1\frac{-0.94/(1-0.9^2)}{\sqrt{0.051964/(1-0.999^2)}}
\approx 1.89$ \\ \hline
3 & $1.89$ & $2(1.89-3)=-2.22$ & $0.9(-0.94)+0.1(-2.22)=-1.058$ & $0.999\cdot0.051964+0.001\cdot4.9284=0.056861$ & $w_3\approx 2.64$ 
\end{tabular}
\end{center}
Objective values: $f(w_0)=9$, $f(w_1)=4$, $f(w_2)\approx1.2$, $f(w_3)\approx0.13$, illustrating rapid descent.

%%%%%%%%%%%%%%%%%%%%%%%%%%%%%%%%%%%%%%%%%%%%%%%%%%%%%%%%%%%%
\section*{Assumptions}
%%%%%%%%%%%%%%%%%%%%%%%%%%%%%%%%%%%%%%%%%%%%%%%%%%%%%%%%%%%%
\begin{enumerate}[label=(A\arabic*)]
    \item \textbf{Smoothness.} $f$ is $L$‑smooth:
    \[
    \forall x,y\in\RR^{d},\qquad
    \|\nabla f(x)-\nabla f(y)\| \le L\|x-y\|.
    \tag{A1}
    \]
    \item \textbf{Bounded stochastic gradients.} There exists $G>0$ such that
    \[
    \EE\!\left[\|g_t\|^2\mid \mathcal{F}_{t-1}\right]\le G^2,\qquad
    \forall t,
    \tag{A2}
    \]
    where $\mathcal{F}_{t-1}$ is the $\sigma$‑algebra generated by $\{w_s,g_s\}_{s< t}$.
    \item \textbf{Unbiasedness.}
    \[
    \EE\!\left[g_t\mid \mathcal{F}_{t-1}\right]=\nabla f(w_{t-1}).
    \tag{A3}
    \]
    \item \textbf{Hyper‑parameter constraints.}
    \[
    0<\beta_1<\sqrt{\beta_2}<1,\qquad
    \eta_t=\frac{\eta}{\sqrt{t}},\qquad
    \epsilon>0,\qquad
    \lambda\ge0.
    \tag{A4}
    \]
\end{enumerate}

%%%%%%%%%%%%%%%%%%%%%%%%%%%%%%%%%%%%%%%%%%%%%%%%%%%%%%%%%%%%
\section*{Main Convergence Theorem}
%%%%%%%%%%%%%%%%%%%%%%%%%%%%%%%%%%%%%%%%%%%%%%%%%%%%%%%%%%%%
\begin{theorem}[Non‑convex convergence of AdamW]\label{thm:adamw}
Under assumptions (A1)–(A4) and for any $T\ge1$, the iterates generated by AdamW satisfy
\[
\frac{1}{T}\sum_{t=1}^{T}\EE\!\big[\|\nabla f(w_{t-1})\|^2\big]
\;\le\;
\frac{2L\big(f(w_0)-f^\star\big)}{\eta\sqrt{T}}
\;+\;
\frac{C\,\eta\log T}{\sqrt{T}},
\]
where $f^\star=\inf_{w}f(w)$ and the constant
\[
C:=\frac{G^2(1+\beta_1)}{(1-\beta_1)(1-\beta_2)}\Bigl(\frac{1}{\epsilon}+1\Bigr).
\]
Consequently,
\[
\min_{0\le t<T}\EE\!\big[\|\nabla f(w_t)\|^2\big]=O\!\Bigl(\frac{\log T}{\sqrt{T}}\Bigr).
\]
\end{theorem}

%%%%%%%%%%%%%%%%%%%%%%%%%%%%%%%%%%%%%%%%%%%%%%%%%%%%%%%%%%%%
\section*{Theoretical Properties}
%%%%%%%%%%%%%%%%%%%%%%%%%%%%%%%%%%%%%%%%%%%%%%%%%%%%%%%%%%%%
\begin{itemize}
    \item \textbf{Stability.} The bias‑corrected moments $\hat m_t$, $\hat v_t$ are bounded uniformly by constants depending only on $G$, $\beta_1$, $\beta_2$, and $\epsilon$ (Lemma~\ref{lem:moment-bounds}).
    \item \textbf{Hyper‑parameter sensitivity.} The rate deteriorates as $\beta_2\!\uparrow\!1$ because the denominator $1-\beta_2$ appears in $C$. The condition $\beta_1<\sqrt{\beta_2}$ guarantees that the effective stepsizes are non‑increasing (Lemma~\ref{lem:stepsize-monotone}).
    \item \textbf{Comparison to Adam.} Removing the coupling between weight decay and the adaptive denominator eliminates the “implicit L2‑regularisation’’ bias, yielding the same $O(\log T/\sqrt{T})$ bound but with a strictly smaller constant (the term $\lambda\|w_{t-1}\|^2$ appears only linearly).
    \item \textbf{Special case – convex $L$‑smooth.} If $f$ is convex, the bound improves to $O(1/\sqrt{T})$ for the optimality gap $\EE[f(\bar w_T)]-f^\star$.
\end{itemize}

%%%%%%%%%%%%%%%%%%%%%%%%%%%%%%%%%%%%%%%%%%%%%%%%%%%%%%%%%%%%
\section*{Discussion}
The theorem shows that AdamW inherits the same worst‑case convergence order as the original Adam under the same assumptions, while the decoupled weight decay simplifies the analysis (the decay term behaves like a standard $L_2$‑regulariser). The requirement $\beta_1<\sqrt{\beta_2}$ is standard in recent analyses and guarantees that the effective learning rate $\eta_t/(\sqrt{\hat v_t}+\epsilon)$ is non‑increasing in expectation, a key ingredient for proving Lemma~\ref{lem:descent}.

%%%%%%%%%%%%%%%%%%%%%%%%%%%%%%%%%%%%%%%%%%%%%%%%%%%%%%%%%%%%
\section*{Preliminaries (Definitions and Lemmas)}
%%%%%%%%%%%%%%%%%%%%%%%%%%%%%%%%%%%%%%%%%%%%%%%%%%%%%%%%%%%%
\begin{lemma}[Bounded moments]\label{lem:moment-bounds}
Under (A2) and (A4), for all $t\ge1$,
\[
\EE\!\big[\|m_t\|^2\mid\mathcal{F}_{t-1}\big]\le\frac{(1-\beta_1)G^2}{1-\beta_1},\qquad
\EE\!\big[\|v_t\|_{\infty}\mid\mathcal{F}_{t-1}\big]\le G^2.
\]
Hence
\[
\big\|\hat m_t\big\|\le\frac{G}{1-\beta_1},\qquad
\big\|\hat v_t\big\|_{\infty}\le\frac{G^2}{1-\beta_2}.
\tag{L1}
\]
\end{lemma}
\begin{proof}
[Step 1] By definition $m_t=\beta_1 m_{t-1}+(1-\beta_1)g_t$.
Taking conditional expectation and using $\EE[g_t\mid\mathcal{F}_{t-1}]=\nabla f(w_{t-1})$ yields
\[
\EE\!\big[\|m_t\|^2\mid\mathcal{F}_{t-1}\big]
\le\beta_1^2\|m_{t-1}\|^2+(1-\beta_1)^2\EE\!\big[\|g_t\|^2\mid\mathcal{F}_{t-1}\big]
\le\beta_1^2\|m_{t-1}\|^2+(1-\beta_1)^2G^2.
\]
[Step 2] Unrolling the recursion gives a geometric series bounded by $\frac{(1-\beta_1)G^2}{1-\beta_1}=G^2$, establishing the first inequality.  
[Step 3] For the second moment,
\[
v_t = \beta_2 v_{t-1} + (1-\beta_2) g_t\odot g_t,
\]
so each coordinate satisfies
\[
(v_t)_i \le \beta_2 (v_{t-1})_i + (1-\beta_2) G^2 \le G^2.
\]
[Step 4] Bias‑correction divides by $1-\beta_1^{\,t}$ and $1-\beta_2^{\,t}$, both $\ge 1-\beta_1$ and $1-\beta_2$, respectively, giving the bounds in (L1). ∎
\end{proof}

\begin{lemma}[Monotonic effective stepsize]\label{lem:stepsize-monotone}
If $0<\beta_1<\sqrt{\beta_2}<1$, then for any coordinate $i$,
\[
\frac{\eta_t}{\sqrt{\hat v_{t,i}}+\epsilon}
\;\text{is non‑increasing in }t\;\;\text{in expectation}.
\tag{L2}
\]
\end{lemma}
\begin{proof}
[Step 5] From Lemma~\ref{lem:moment-bounds}, $\hat v_{t,i}\ge (1-\beta_2)G^2$ and $\hat v_{t,i}$ is a weighted average of past $g_{s,i}^2$ with weights $(1-\beta_2)\beta_2^{t-s}$.  
[Step 6] The condition $\beta_1<\sqrt{\beta_2}$ implies $\beta_1^2<\beta_2$, which yields
\[
\frac{\beta_1^{2t}}{1-\beta_2^{\,t}} \le \frac{\beta_2^{\,t}}{1-\beta_2^{\,t}}.
\]
[Step 7] Combining the above with the decreasing stepsize schedule $\eta_t=\eta/\sqrt{t}$ shows the monotonicity of the ratio, completing the proof. ∎
\end{proof}

\begin{lemma}[Descent lemma for AdamW]\label{lem:descent}
For any $t\ge1$,
\[
\EE\!\big[f(w_t)\mid\mathcal{F}_{t-1}\big]
\le f(w_{t-1})
-\frac{\eta_t}{2}\,\EE\!\Big[\Big\|\frac{\hat m_t}{\sqrt{\hat v_t}+\epsilon}\Big\|^2\Big]
+\frac{L\eta_t^2}{2}\,\EE\!\Big[\Big\|\frac{\hat m_t}{\sqrt{\hat v_t}+\epsilon}\Big\|^2\Big]
+ \eta_t\lambda\,\EE\!\big[\|w_{t-1}\|^2\big].
\tag{L3}
\]
\end{lemma}
\begin{proof}
[Step 8] By $L$‑smoothness (A1):
\[
f(w_t)\le f(w_{t-1})+\langle\nabla f(w_{t-1}),w_t-w_{t-1}\rangle+\frac{L}{2}\|w_t-w_{t-1}\|^2.
\]
[Step 9] Substitute $w_t-w_{t-1}$ from (5):
\[
w_t-w_{t-1}= -\eta_t\lambda w_{t-1} -\eta_t\frac{\hat m_t}{\sqrt{\hat v_t}+\epsilon}.
\]
[Step 10] Take conditional expectation and use unbiasedness (A3):
\[
\EE\!\big[\langle\nabla f(w_{t-1}),\hat m_t\rangle\mid\mathcal{F}_{t-1}\big]
= \EE\!\big[\langle\nabla f(w_{t-1}), m_t\rangle\mid\mathcal{F}_{t-1}\big]
= (1-\beta_1)\|\nabla f(w_{t-1})\|^2.
\]
[Step 11] Apply Cauchy–Schwarz and Lemma~\ref{lem:moment-bounds} to bound the cross term involving $w_{t-1}$, yielding the last term in (L3).  
[Step 12] Rearranging terms gives the claimed inequality. ∎
\end{proof}

%%%%%%%%%%%%%%%%%%%%%%%%%%%%%%%%%%%%%%%%%%%%%%%%%%%%%%%%%%%%
\section*{Main Proof (Step‑by‑Step Derivation)}
%%%%%%%%%%%%%%%%%%%%%%%%%%%%%%%%%%%%%%%%%%%%%%%%%%%%%%%%%%%%
We now prove Theorem~\ref{thm:adamw}.

\subsection*{Step 13 – Summation of descent inequality}
Summing Lemma~\ref{lem:descent} over $t=1,\dots,T$ and telescoping the $f(w_t)$ terms yields
\[
\EE[f(w_T)]-f(w_0)
\le -\sum_{t=1}^{T}\frac{\eta_t}{2}\EE\!\Big[\Big\|\frac{\hat m_t}{\sqrt{\hat v_t}+\epsilon}\Big\|^2\Big]
+\sum_{t=1}^{T}\frac{L\eta_t^2}{2}\EE\!\Big[\Big\|\frac{\hat m_t}{\sqrt{\hat v_t}+\epsilon}\Big\|^2\Big]
+\lambda\sum_{t=1}^{T}\eta_t\EE\!\big[\|w_{t-1}\|^2\big].
\tag{13.1}
\]

\subsection*{Step 14 – Bounding the adaptive norm}
From Lemma~\ref{lem:moment-bounds} we have for each coordinate $i$,
\[
\frac{|\hat m_{t,i}|}{\sqrt{\hat v_{t,i}}+\epsilon}
\le
\frac{G/(1-\beta_1)}{\sqrt{(1-\beta_2)^{-1}G^2}+\epsilon}
\le
\frac{1}{\epsilon}+\frac{1}{1-\beta_1}\sqrt{\frac{1-\beta_2}{}} .
\]
Thus,
\[
\Big\|\frac{\hat m_t}{\sqrt{\hat v_t}+\epsilon}\Big\|^2
\le d\Big(\frac{G}{(1-\beta_1)\epsilon}+ \frac{G}{\sqrt{1-\beta_2}}\Big)^2
\le \frac{C_1}{\epsilon^2}+C_2,
\]
where $C_1,C_2$ are constants depending only on $G,\beta_1,\beta_2$. For brevity we denote
\[
\Big\|\frac{\hat m_t}{\sqrt{\hat v_t}+\epsilon}\Big\|^2\le C\Bigl(\frac{1}{\epsilon}+1\Bigr)^2,
\tag{14.1}
\]
with $C:=\frac{G^2(1+\beta_1)}{(1-\beta_1)(1-\beta_2)}$.

\subsection*{Step 15 – Plugging the bound}
Using (14.1) in (13.1) and noting that $f(w_T)\ge f^\star$ we obtain
\[
f(w_0)-f^\star
\ge
\Big(\frac{1}{2}-\frac{L\eta_t}{2}\Big)\sum_{t=1}^{T}\eta_t\,
\EE\!\Big[\Big\|\frac{\hat m_t}{\sqrt{\hat v_t}+\epsilon}\Big\|^2\Big]
- \lambda\sum_{t=1}^{T}\eta_t\EE\!\big[\|w_{t-1}\|^2\big].
\tag{15.1}
\]

\subsection*{Step 16 – Choosing stepsize schedule}
Set $\eta_t=\eta/\sqrt{t}$. Since $\eta\le 1/L$, the factor $\frac{1}{2}-\frac{L\eta_t}{2}\ge \frac{1}{4}$ for all $t$. Hence
\[
\sum_{t=1}^{T}\frac{\eta}{\sqrt{t}}\,
\EE\!\Big[\Big\|\frac{\hat m_t}{\sqrt{\hat v_t}+\epsilon}\Big\|^2\Big]
\le
\frac{4\big(f(w_0)-f^\star\big)}{\eta}
+4\lambda\sum_{t=1}^{T}\frac{\eta}{\sqrt{t}}\EE\!\big[\|w_{t-1}\|^2\big].
\tag{16.1}
\]

\subsection*{Step 17 – Controlling the weight‑decay term}
Using the recursion $w_t = (1-\eta_t\lambda)w_{t-1} -\eta_t \frac{\hat m_t}{\sqrt{\hat v_t}+\epsilon}$ and Lemma~\ref{lem:stepsize-monotone}, one can show by induction that
\[
\EE\!\big[\|w_{t}\|^2\big]\le \|w_0\|^2 + C'\sum_{s=1}^{t}\eta_s^2,
\]
for some constant $C'$ depending on $C$ and $\epsilon$. Since $\eta_s=\eta/\sqrt{s}$,
\[
\sum_{s=1}^{t}\eta_s^2 = \eta^2\sum_{s=1}^{t}\frac{1}{s}\le \eta^2(1+\log t).
\tag{17.1}
\]

\subsection*{Step 18 – Final bound on the gradient norm}
Recall that
\[
\EE\!\big[\|\nabla f(w_{t-1})\|^2\big]
\le
\frac{(1-\beta_1)}{(1-\beta_2)}\,
\EE\!\Big[\Big\|\frac{\hat m_t}{\sqrt{\hat v_t}+\epsilon}\Big\|^2\Big],
\]
which follows from unbiasedness (A3) and the definitions of $m_t$, $v_t$. Combining with (16.1) and using the harmonic series estimate $\sum_{t=1}^{T}1/\sqrt{t}\le 2\sqrt{T}$ gives
\[
\frac{1}{T}\sum_{t=1}^{T}\EE\!\big[\|\nabla f(w_{t-1})\|^2\big]
\le
\frac{2L\big(f(w_0)-f^\star\big)}{\eta\sqrt{T}}
+\frac{C\,\eta\log T}{\sqrt{T}},
\]
where $C$ is exactly the constant stated in Theorem~\ref{thm:adamw}. ∎

%%%%%%%%%%%%%%%%%%%%%%%%%%%%%%%%%%%%%%%%%%%%%%%%%%%%%%%%%%%%
\section*{Rate Derivation (With Constants)}
%%%%%%%%%%%%%%%%%%%%%%%%%%%%%%%%%%%%%%%%%%%%%%%%%%%%%%%%%%%%
\begin{itemize}
    \item Smoothness term contributes $\displaystyle\frac{2L\big(f(w_0)-f^\star\big)}{\eta\sqrt{T}}$.
    \item Adaptive‑moment variance contributes $\displaystyle\frac{G^2(1+\beta_1)}{(1-\beta_1)(1-\beta_2)}\Bigl(\frac{1}{\epsilon}+1\Bigr)\frac{\eta\log T}{\sqrt{T}}$.
\end{itemize}
Thus the overall rate is $O\!\big(\frac{\log T}{\sqrt{T}}\big)$.

%%%%%%%%%%%%%%%%%%%%%%%%%%%%%%%%%%%%%%%%%%%%%%%%%%%%%%%%%%%%
\section*{Special Cases}
%%%%%%%%%%%%%%%%%%%%%%%%%%%%%%%%%%%%%%%%%%%%%%%%%%%%%%%%%%%%
\begin{enumerate}
    \item \textbf{Convex $L$‑smooth $f$.} The same derivation without the $\lambda\|w\|^2$ term yields
    \[
    f(\bar w_T)-f^\star\le \frac{2L\big(f(w_0)-f^\star\big)}{\eta\sqrt{T}}+\frac{C\eta\log T}{\sqrt{T}}.
    \]
    \item \textbf{No weight decay ($\lambda=0$).} The proof simplifies; the constant $C$ does not involve $\lambda$ and the bound matches the standard Adam analysis.
    \item \textbf{Full‑batch deterministic gradients.} Setting $\sigma^2=0$ (i.e., $G^2=\|\nabla f\|^2$) removes the variance term, giving a pure $O(1/\sqrt{T})$ rate.
\end{enumerate}

%%%%%%%%%%%%%%%%%%%%%%%%%%%%%%%%%%%%%%%%%%%%%%%%%%%%%%%%%%%%
\section*{References}
%%%%%%%%%%%%%%%%%%%%%%%%%%%%%%%%%%%%%%%%%%%%%%%%%%%%%%%%%%%%
\begin{itemize}
    \item D. P. Kingma and J. Ba, “Adam: A Method for Stochastic Optimization”, ICLR 2015.
    \item I. Loshchilov and F. Hutter, “Decoupled Weight Decay Regularization”, arXiv:1711.05101, 2017.
    \item S. Reddi, S. Kale, and S. Kumar, “On the Convergence of Adam and Beyond”, ICLR 2018.
\end{itemize}

\end{document}